\subsection{Mathematical Model}
\label{sec:model}

We relate acoustic measurements to supplied valence and arousal scores. We fit a separate model for each response because a recording can have a negative emotional character and high activation at the same time. Ordinary least squares is our starting point, and ridge regression adds a penalty that can make the fitted coefficients less sensitive to related predictors.

The completed pilot uses the earlier 92-recording selection. The framework also applies to the revised question across genres, but the broader sample still needs to be prepared and evaluated. A sequence-model extension is proposed below.

\subsubsection{Predictors and responses}
Let $n$ denote the number of eligible excerpts in the finalized sample, and let $\bm{x}_i\in\mathbb{R}^7$ contain the measurements listed in Table~\ref{tab:predictors}. These seven features were extracted in the preliminary analysis; final inclusion depends on audio availability, label checks, and review of repeated recordings. No mood tags, quadrant labels, or other emotion-derived metadata will be used as predictors. The released numerical labels we inspected lie on a $[0,1]$ scale. We will preserve that scale and document it alongside the annotation procedure.

\begin{table}[ht]
\centering
\caption{Initial acoustic predictors, computed from each excerpt.}
\label{tab:predictors}
\begin{tabular}{ll}
\toprule
Measurement & Musical interpretation \\
\midrule
Mean frame RMS amplitude & Typical signal amplitude \\
Standard deviation of frame RMS & Variation in signal amplitude \\
Mean spectral centroid (Hz) & Spectral brightness \\
Power fraction in $20\leq f<250$ Hz & Bass-band content \\
Mean positive spectral flux & Changes in spectral magnitude \\
Detected onsets per second & Rate of detected sound attacks \\
Estimated tempo (BPM) & Estimated beat rate \\
\bottomrule
\end{tabular}
\end{table}

RMS amplitude describes signal amplitude and serves as an approximate loudness-related measurement. Estimated tempo is subject to beat-tracking errors. Bass-band content measures the proportion of spectrogram power in the stated frequency range. These quantities describe aspects of the audio; arousal remains a response label. Harmony and mode are possible later additions, but we have not yet validated features that would support claims about minor chords.

\subsubsection{Defining the audio measurements}

The preliminary extraction converts audio to mono at $22{,}050$ samples per second. It uses a short-time Fourier transform with a Hann window of length $N=2{,}048$ and a hop length of $512$ samples. Let $S_{kt}\geq0$ denote the resulting magnitude at frequency bin $k$ and frame $t$, with frequency $f_k$. Let $L$ be the number of frames and $T$ the excerpt duration in seconds.

For the windowed samples $u_{t\ell}$ in a frame, its RMS amplitude is
\begin{equation}
 R_t=\sqrt{\frac{1}{N}\sum_{\ell=0}^{N-1}u_{t\ell}^2}.
\end{equation}
The extraction computes the equivalent quantity from the spectrogram. Its two summaries are
\begin{equation}
 \overline R=\frac{1}{L}\sum_{t=1}^{L}R_t,
 \qquad
 s_R=\sqrt{\frac{1}{L}\sum_{t=1}^{L}(R_t-\overline R)^2}.
\end{equation}
Here $s_R$ is a population standard deviation across frames. Both measurements depend on the recording's amplitude scale and production, which matters when comparing mastered tracks.

The spectral centroid for a non-silent frame is
\begin{equation}
 C_t=\frac{\sum_k f_k S_{kt}}{\sum_k S_{kt}},
 \qquad \overline C=\frac{1}{L}\sum_{t=1}^{L}C_t.
\end{equation}
Silent frames are assigned centroid zero in the extraction. Larger values indicate that spectral magnitude is concentrated at higher frequencies. The bass-band fraction is
\begin{equation}
 B=\frac{\sum_t\sum_{k:\,20\leq f_k<250}S_{kt}^2}
 {\sum_t\sum_k S_{kt}^2}.
\end{equation}
The implementation adds a small constant to the denominator to handle near-silence. Positive spectral flux is summarized by
\begin{equation}
 F=\frac{1}{L-1}\sum_{t=2}^{L}\sum_k
 \max\{S_{kt}-S_{k,t-1},0\}.
\end{equation}
This particular flux measure uses unnormalized spectral magnitudes, so amplitude changes can affect it. We will inspect its correlation with the RMS features before interpreting its coefficient.

If the onset detector identifies $D$ attacks, the onset-rate feature is $D/T$. The tempo estimate is the beat-rate estimate returned by the audio analysis algorithm, measured in BPM. Onset rate and tempo describe different aspects of rhythm: a recording can contain several detected attacks within one beat. We will inspect beat-tracking errors, including estimates at half or twice a plausible tempo. All case-study excerpts will use this same feature-extraction procedure.

For training-set size $m$, define the standardized predictor
\begin{equation}
 z_{ij}=\frac{x_{ij}-\mu_j}{s_j},
 \qquad j=1,\ldots,p,
\end{equation}
where $\mu_j=m^{-1}\sum_{i=1}^m x_{ij}$ and $s_j=\{m^{-1}\sum_{i=1}^m(x_{ij}-\mu_j)^2\}^{1/2}$ are the training mean and standard deviation. Constant predictors are removed, so $p\leq7$. The same training values are used to transform held-out songs. A coefficient $\beta_{jr}$ therefore represents the fitted change in response $r$ associated with a one-training-standard-deviation increase in predictor $j$, holding the other standardized predictors fixed.

\subsubsection{Ordinary least squares}

For either response, let $y_i$ denote the supplied score for recording $i$. Our initial model is
\begin{equation}
y_i=\beta_0+\sum_{j=1}^{p}\beta_jz_{ij}+\varepsilon_i.
\end{equation}
The intercept describes the prediction at the training-average feature values. Each coefficient describes the change in prediction associated with a one-standard-deviation increase in a feature, with the other included features fixed. The residual captures variation that the model does not explain.

Ordinary least squares chooses coefficients to minimize
\begin{equation}
\frac{1}{m}\sum_{i=1}^{m}
\left(y_i-\beta_0-\sum_{j=1}^{p}\beta_jz_{ij}\right)^2.
\end{equation}
This gives larger errors more weight and prevents positive and negative errors from canceling. With centered response vector $\bm y_c$ and standardized feature matrix $Z$, the coefficient equations are
\[
Z^{\mathsf T}Z\widehat{\bm\beta}=Z^{\mathsf T}\bm y_c.
\]
If the predictors are linearly independent, their solution is
\[
\widehat{\bm\beta}=(Z^{\mathsf T}Z)^{-1}Z^{\mathsf T}\bm y_c.
\]
The implementation uses a numerical least-squares solver rather than explicitly computing this inverse. Correlated measurements can make the coefficients unstable, which motivates the ridge model.

\subsubsection{Ridge estimation}

Ridge regression adds a squared-coefficient penalty to least squares; van Wieringen reviews this framework and its properties \cite{vanwieringen2026ridge}. Below we derive the estimator using our chosen normalization of the loss.
For response index $r\in\{v,a\}$, define $y_{iv}=v_i$ and $y_{ia}=a_i$. The model is
\begin{equation}
 y_{ir}=\beta_{0r}+\bm{z}_i^{\mathsf T}\bm{\beta}_r+\varepsilon_{ir}.
 \label{eq:linear}
\end{equation}
We estimated the coefficients by minimizing
\begin{equation}
 (\widehat\beta_{0r},\widehat{\bm\beta}_r)
 =\operatorname*{arg\,min}_{b,\bm\beta}
 \left\{\frac{1}{m}\sum_{i=1}^{m}
 (y_{ir}-b-\bm z_i^{\mathsf T}\bm\beta)^2
 +\lambda_r\lVert\bm\beta\rVert_2^2\right\},\qquad\lambda_r>0.
 \label{eq:ridge}
\end{equation}
The intercept is unpenalized. The penalty limits coefficient size, which is useful because measurements such as amplitude and onset activity may be correlated. Equation~\eqref{eq:linear} is a linear approximation to the relationship between the chosen measurements and the supplied labels. The residual term includes both variation not captured by these predictors and error in the measurements or labels.

\begin{prop}
For any $\lambda_r>0$, problem~\eqref{eq:ridge} has a unique solution, even if the standardized predictors are linearly dependent.
\end{prop}
\begin{proof}
Let $Z$ be the $m\times p$ training matrix with centered columns and let $\bm y_{c,r}=\bm y_r-\bar y_r\bm 1$. The derivative with respect to the intercept is $-2m^{-1}\sum_i(y_{ir}-b-\bm z_i^{\mathsf T}\bm\beta)$. Since $\sum_i\bm z_i=\bm0$, setting it to zero gives $\widehat\beta_{0r}=\bar y_r$. The normal equations for the remaining coefficients are
\[
 (Z^{\mathsf T}Z+m\lambda_r I_p)\widehat{\bm\beta}_r
 =Z^{\mathsf T}\bm y_{c,r}.
\]
For every nonzero $\bm u\in\mathbb R^p$,
\[
 \bm u^{\mathsf T}(Z^{\mathsf T}Z+m\lambda_r I_p)\bm u
 =\lVert Z\bm u\rVert_2^2+m\lambda_r\lVert\bm u\rVert_2^2>0.
\]
The matrix is positive definite and invertible. Thus
\[
 \widehat{\bm\beta}_r
 =(Z^{\mathsf T}Z+m\lambda_r I_p)^{-1}Z^{\mathsf T}\bm y_{c,r},
\]
and both the coefficient vector and intercept are unique.
\end{proof}

The penalty shrinks feature coefficients toward zero. A larger penalty brings predictions closer to the training-mean intercept. This can reduce sensitivity to the training sample, although too much shrinkage can remove useful information. We select the penalty through validation. With the average-error convention in Equation~\eqref{eq:ridge}, software using a sum-of-errors objective requires $\alpha=m\lambda_r$.

\subsubsection{A proposed LSTM extension}

Excerpt-level summaries do not preserve the order of acoustic changes. A recording that gradually becomes louder can have a similar mean amplitude to one that starts loudly and becomes quiet. To examine whether that order improves prediction, we propose a long short-term memory network, or LSTM \cite{hochreiter1997lstm}.

Each recording would be divided into consecutive windows, with a feature vector $\bm z_{it}$ for window $t$. The model updates a hidden representation and memory state:
\[
(\bm h_{it},\bm c_{it})
=
\operatorname{LSTM}_{\theta}
(\bm z_{it},\bm h_{i,t-1},\bm c_{i,t-1}).
\]
These states allow earlier measurements to influence the final representation. Two output layers would then predict the overall responses:
\[
\widehat v_i=b_v+\bm w_v^{\mathsf T}\bm h_{iT},
\qquad
\widehat a_i=b_a+\bm w_a^{\mathsf T}\bm h_{iT}.
\]
Training would minimize the average squared error across both responses. The model has not yet been trained. All windows from one recording will remain in the same evaluation group, and preprocessing will use training data only. With static labels, this tests prediction of overall emotion scores; it does not validate emotion predictions for individual windows.

\subsubsection{Evaluation and interpretation}
The completed pilot used five outer folds and three inner folds, grouping observations by complete artist-credit strings after normalizing whitespace and letter case. The pilot contains 63 such groups. We will rebuild the groups for the broader sample and use the same evaluation design where sample size permits. No artist-credit group appears in both training and validation or test data within a split. The primary outer split uses random seed 195. Inner grouped validation selected $\lambda_r$ from $\{10^{-4},10^{-3},\ldots,10^4\}$ by pooled validation mean squared error. Feature means and scales were fitted separately within each training fold. All predictors were finite, so no imputation was required. Aliases and overlapping personnel across different artist credits were not resolved, which limits the strength of the artist-disjoint interpretation. The personal case studies have not yet been completed and were not used in pilot model selection.

For held-out predictions, we report
\begin{equation}
 \mathrm{RMSE}_r=
 \sqrt{\frac{1}{n_{\mathrm{test}}}\sum_{i\in\mathrm{test}}
 (y_{ir}-\widehat y_{ir})^2}.
\end{equation}
We compared this error with a baseline that predicts the outer training mean. Ridge predictions can fall outside $[0,1]$; we report their frequency and evaluate the raw predictions consistently. Coefficients are summarized across outer folds to show whether their signs and magnitudes are stable. To assess predictive usefulness, we also refitted and retuned models with each feature family omitted, using the same outer folds. Correlated predictors can substitute for one another, so a small error change does not establish that a musical characteristic is irrelevant.

We grouped the predictors into amplitude ($\overline R,s_R$), rhythm (onset rate and tempo), spectral change and brightness ($F,\overline C$), and bass-band content ($B$). For family $g$, define the held-out error change
\begin{equation}
 \Delta_{g,r}=\mathrm{RMSE}_{r,-g}-\mathrm{RMSE}_{r,\mathrm{full}}.
\end{equation}
A positive value indicates that omitting the family worsens prediction under the stated evaluation procedure. We report the changes across folds as well as the aggregate value. A negative value is possible if the omitted features add noise or do not generalize. These comparisons will complement coefficient summaries rather than provide a ranking of causal effects.

We also report out-of-fold $R^2$:
\begin{equation}
 R^2_r=1-\frac{\sum_i(y_{ir}-\widehat y^{\mathrm{OOF}}_{ir})^2}
 {\sum_i(y_{ir}-\bar y_r)^2},
\end{equation}
provided the response has nonzero variation. Here each prediction is obtained from an outer model that did not train on that example, and $\bar y_r$ is the mean of the evaluated responses. This statistic can be negative and is not a classification accuracy. The training-mean baseline will still be evaluated separately in every outer fold, since it is the deployable constant predictor for that fold.

\subsubsection{Case-study predictions}
After pilot evaluation, each primary model was refitted on the full provisional sample, with five-fold grouped validation selecting its penalty. These pilot refits are saved for exploratory case-study predictions; their training errors are not used to evaluate predictive performance. Case studies for the revised question will use models refitted on the broader sample after evaluation. The case-study recordings will remain excluded from all training and tuning. For a case-study excerpt $\ast$, define $c_{jr}=\widehat\beta_{jr}z_{\ast j}$. Then
\begin{equation}
 \widehat y_{\ast r}-\widehat\beta_{0r}
 =\sum_{j=1}^{p}c_{jr}.
 \label{eq:contributions}
\end{equation}
This identity gives an exact decomposition of the fitted linear prediction relative to its training-mean intercept. A negative valence contribution indicates that a measured feature lowers the model's prediction for that excerpt. The contribution describes an association within the fitted model. Establishing a causal effect of changing the musical characteristic would require an experimental design beyond the scope of this study. Each case study will use a prespecified excerpt of approximately 30 seconds, identified by its timestamp before model predictions are inspected. We will check whether its standardized features lie within the observed training ranges and compare the predicted valence and arousal with an independently written listening description. We will report the intercept, individual feature contributions, and residual difference if a comparable emotion annotation is available. An informal description alone will not be treated as a numerical ground-truth score. Prediction errors will remain part of the case-study analysis and will be discussed in relation to the coverage of the feature set and the training sample.

\subsubsection{Temporal descriptors and a section-based extension}

To investigate acoustic variation within the excerpts, we divided each recording into ten equal-duration bins. Let $r_{i\ell}$, $c_{i\ell}$, $b_{i\ell}$, and $o_{i\ell}$ denote mean RMS amplitude, mean centroid, bass-band power fraction, and onset rate in bin $\ell$, respectively. For any bin-level measurement $h_{i\ell}$, define its time slope over normalized bin centers $t_\ell=(\ell-1/2)/10$ as
\begin{equation}
 s_i(h)=\frac{\sum_{\ell=1}^{10}(t_\ell-\bar t)(h_{i\ell}-\bar h_i)}
 {\sum_{\ell=1}^{10}(t_\ell-\bar t)^2}.
\end{equation}
We used slopes of amplitude normalized by excerpt mean amplitude, centroid, and onset rate. We also measured
\begin{equation}
 J_i=\max_{\ell=1,\ldots,9}
 20\log_{10}\!\left(\frac{r_{i,\ell+1}+10^{-8}}{r_{i\ell}+10^{-8}}\right),
 \qquad
 K_i=\max_{\ell=1,\ldots,9}(b_{i,\ell+1}-b_{i\ell}).
\end{equation}
These five descriptors were added to the seven original features in an exploratory ridge model. The extension was introduced after examining the primary results and is reported separately. It used the same outer groups and nested tuning procedure. No new emotion responses were created for the bins.

A verse, chorus, bridge, buildup, drop, or breakdown is a musical interpretation of a section. The appropriate labels depend on the recording. The descriptors above measure patterns that can motivate section review, but they do not assign verified section labels. For later annotated full-song examples, we define a section contrast for feature $j$ by
\begin{equation}
 d_j=x_{\mathrm{drop},j}-x_{\mathrm{preceding\ buildup},j}.
\end{equation}
Using the same fitted primary model and scaling, the difference between the two section-input predictions is exactly
\begin{equation}
 \widehat y_{\mathrm{drop},r}-\widehat y_{\mathrm{buildup},r}
 =\sum_j\widehat\beta_{jr}\frac{d_j}{s_j}.
\end{equation}
This expresses how changes in measured sound affect the model's output. The model was trained on static excerpt labels, so the difference cannot be interpreted as a validated change in listener emotion without section-level annotations. Shorter sections also introduce a distribution shift relative to the approximately 30-second training excerpts.

